% --- Archon physics formalization source begin ---
% archon:physics
% archon:covers PhyXMiniProblems/problem_phyx_mini_0578.lean
% archon:source-report reports/phyx_mini/problem_phyx_mini_0578.source.json

\chapter{Physics problem phyx\_mini\_0578}
\label{ch:PhyXMiniProblems_problem_phyx_mini_0578}

\paragraph{Problem source.}
\#\# Physical scenario

Excited sodium atoms emit two closely spaced spectrum lines called the sodium doublet with wavelengths \$588.995 \textbackslash{}, \textbackslash{}text\{nm\}\$ and \$589.592 \textbackslash{}, \textbackslash{}text\{nm\}\$.

\#\# Question

This energy difference occurs because the electron's spin magnetic moment can be oriented either parallel or antiparallel to the internal magnetic field associated with the electron's orbital motion.

\#\# Answer choices

- A: A: 9T

- B: B: 1.8T

- C: C: 18T

- D: D: 36T

\#\# Figure caption (auxiliary; use the image as primary evidence)

The image depicts an energy level diagram commonly used in atomic physics:

1. **Energy Levels**:
   - There are two main horizontal lines representing energy levels.
   - The top energy level is labeled with \textbackslash{}( n = 3, \textbackslash{}ell = 1 \textbackslash{}).
   - The bottom energy level is labeled with \textbackslash{}( n = 3, \textbackslash{}ell = 0 \textbackslash{}).

2. **Transitions**:
   - Two vertical lines with downward arrows indicate transitions between these energy levels.
   - Each arrow represents the emission of a photon as an electron transitions to a lower energy state.

3. **Wavelengths**:
   - The left transition has a wavelength labeled as \textbackslash{}( 588.995 \textbackslash{}, \textbackslash{}text\{nm\} \textbackslash{}).
   - The right transition has a wavelength labeled as \textbackslash{}( 589.592 \textbackslash{}, \textbackslash{}text\{nm\} \textbackslash{}).

4. **Bracket**:
   - A curly bracket on the left groups the top two lines, indicating they belong to the same energy level.

These features depict how electrons transition between energy states, emitting photons with specific wavelengths.

\#\# Dataset metadata

Quantum Phenomena; Physical Model Grounding Reasoning

\paragraph{Recorded answer/context.}
C: C: 18T

\paragraph{Figure/image path.}
/root/proposal\_for\_physic/science-mango/phyx\_mini\_run/phyx\_data/test\_image/578.png

\paragraph{Formalization target.}
create a compiling Lean file with sorry bodies at `PhyXMiniProblems/problem\_phyx\_mini\_0578.lean`. The Lean declarations must preserve the physical quantities, dimensions or dimensional roles, figure labels, governing-law hypotheses, and final relation expressed by this problem.
Use Mathlib/Physlib names found through LeanExplore where available. If a physics API is missing, introduce faithful local abstractions rather than scalar placeholder aliases.

\begin{theorem}[Physics formalization target]
\label{thm:physics:phyx_mini_0578:target}
\lean{PhyXMiniProblems.ProblemPhyXMini0578.problem_phyx_mini_0578}
\uses{def:physics:phyx-mini-0578:phyxminiproblems-problemphyxmini0578-magneticfieldstrengthinteslas, def:physics:phyx-mini-0578:phyxminiproblems-problemphyxmini0578-sodiumdoubletsetup, def:physics:phyx-mini-0578:phyxminiproblems-problemphyxmini0578-matchesexcitedsodiumscenario, def:physics:phyx-mini-0578:phyxminiproblems-problemphyxmini0578-matchessuppliedsodiumdoubletfigure, def:physics:phyx-mini-0578:phyxminiproblems-problemphyxmini0578-matchessodiumdoubletwavelengthreadouts, def:physics:phyx-mini-0578:phyxminiproblems-problemphyxmini0578-hasphysicalsodiumdoubletquantities, def:physics:phyx-mini-0578:phyxminiproblems-problemphyxmini0578-usesstandardspectroscopicconstants, def:physics:phyx-mini-0578:phyxminiproblems-problemphyxmini0578-satisfiessodiumradiativetransitionlaws, def:physics:phyx-mini-0578:phyxminiproblems-problemphyxmini0578-satisfiesinternalspinzeemanlaw, def:physics:phyx-mini-0578:phyxminiproblems-problemphyxmini0578-inferredinternalfieldinteslas, def:physics:phyx-mini-0578:phyxminiproblems-problemphyxmini0578-isuniquematchingdisplayedfield}
The assigned autoformalize agent should translate this physics problem into Lean declarations in the covered file, with theorem and lemma proof bodies written as `by sorry`.
\end{theorem}
\begin{proof}
This is an autoformalization task, not a proof task. Produce statements that can later be proved by the physics prover without weakening the physical model.
\end{proof}

% --- Archon declaration topology begin ---
\section{Declaration topology}
\begin{definition}[energy Dimension]
  \label{def:physics:phyx-mini-0578:phyxminiproblems-problemphyxmini0578-energydimension}
  \lean{PhyXMiniProblems.ProblemPhyXMini0578.energyDimension}
  The physical dimension of energy, M L² T⁻²
\end{definition}
\begin{proof}
  This is the declared model definition of energy Dimension.
\end{proof}
\begin{definition}[action Dimension]
  \label{def:physics:phyx-mini-0578:phyxminiproblems-problemphyxmini0578-actiondimension}
  \lean{PhyXMiniProblems.ProblemPhyXMini0578.actionDimension}
  \uses{def:physics:phyx-mini-0578:phyxminiproblems-problemphyxmini0578-energydimension}
  The physical dimension of action, energy multiplied by time
\end{definition}
\begin{proof}
  This is the declared model definition of action Dimension.
\end{proof}
\begin{definition}[magnetic Field Strength Dimension]
  \label{def:physics:phyx-mini-0578:phyxminiproblems-problemphyxmini0578-magneticfieldstrengthdimension}
  \lean{PhyXMiniProblems.ProblemPhyXMini0578.magneticFieldStrengthDimension}
  The physical dimension M T⁻¹ C⁻¹ of magnetic-field strength
\end{definition}
\begin{proof}
  This is the declared model definition of magnetic Field Strength Dimension.
\end{proof}
\begin{definition}[magnetic Moment Dimension]
  \label{def:physics:phyx-mini-0578:phyxminiproblems-problemphyxmini0578-magneticmomentdimension}
  \lean{PhyXMiniProblems.ProblemPhyXMini0578.magneticMomentDimension}
  \uses{def:physics:phyx-mini-0578:phyxminiproblems-problemphyxmini0578-energydimension, def:physics:phyx-mini-0578:phyxminiproblems-problemphyxmini0578-magneticfieldstrengthdimension}
  The physical dimension of magnetic moment, energy divided by field
\end{definition}
\begin{proof}
  This is the declared model definition of magnetic Moment Dimension.
\end{proof}
\begin{definition}[Wavelength Quantity]
  \label{def:physics:phyx-mini-0578:phyxminiproblems-problemphyxmini0578-wavelengthquantity}
  \lean{PhyXMiniProblems.ProblemPhyXMini0578.WavelengthQuantity}
  A nonnegative, unit-independent physical wavelength
\end{definition}
\begin{proof}
  This is the declared model definition of Wavelength Quantity.
\end{proof}
\begin{definition}[Planck Constant Quantity]
  \label{def:physics:phyx-mini-0578:phyxminiproblems-problemphyxmini0578-planckconstantquantity}
  \lean{PhyXMiniProblems.ProblemPhyXMini0578.PlanckConstantQuantity}
  \uses{def:physics:phyx-mini-0578:phyxminiproblems-problemphyxmini0578-actiondimension}
  A nonnegative physical quantity with the dimension of Planck's constant
\end{definition}
\begin{proof}
  This is the declared model definition of Planck Constant Quantity.
\end{proof}
\begin{definition}[Magnetic Field Strength Quantity]
  \label{def:physics:phyx-mini-0578:phyxminiproblems-problemphyxmini0578-magneticfieldstrengthquantity}
  \lean{PhyXMiniProblems.ProblemPhyXMini0578.MagneticFieldStrengthQuantity}
  \uses{def:physics:phyx-mini-0578:phyxminiproblems-problemphyxmini0578-magneticfieldstrengthdimension}
  A nonnegative magnitude of magnetic-field strength
\end{definition}
\begin{proof}
  This is the declared model definition of Magnetic Field Strength Quantity.
\end{proof}
\begin{definition}[Magnetic Moment Magnitude Quantity]
  \label{def:physics:phyx-mini-0578:phyxminiproblems-problemphyxmini0578-magneticmomentmagnitudequantity}
  \lean{PhyXMiniProblems.ProblemPhyXMini0578.MagneticMomentMagnitudeQuantity}
  \uses{def:physics:phyx-mini-0578:phyxminiproblems-problemphyxmini0578-magneticmomentdimension}
  A nonnegative magnitude of an electron spin magnetic moment
\end{definition}
\begin{proof}
  This is the declared model definition of Magnetic Moment Magnitude Quantity.
\end{proof}
\begin{definition}[wavelength Readout]
  \label{def:physics:phyx-mini-0578:phyxminiproblems-problemphyxmini0578-wavelengthreadout}
  \lean{PhyXMiniProblems.ProblemPhyXMini0578.wavelengthReadout}
  \uses{def:physics:phyx-mini-0578:phyxminiproblems-problemphyxmini0578-wavelengthquantity}
  Read a physical wavelength in a selected length unit
\end{definition}
\begin{proof}
  This is the declared model definition of wavelength Readout.
\end{proof}
\begin{definition}[wavelength In Meters]
  \label{def:physics:phyx-mini-0578:phyxminiproblems-problemphyxmini0578-wavelengthinmeters}
  \lean{PhyXMiniProblems.ProblemPhyXMini0578.wavelengthInMeters}
  \uses{def:physics:phyx-mini-0578:phyxminiproblems-problemphyxmini0578-wavelengthquantity, def:physics:phyx-mini-0578:phyxminiproblems-problemphyxmini0578-wavelengthreadout}
  Read a physical wavelength in metres
\end{definition}
\begin{proof}
  This is the declared model definition of wavelength In Meters.
\end{proof}
\begin{definition}[wavelength In Nanometers]
  \label{def:physics:phyx-mini-0578:phyxminiproblems-problemphyxmini0578-wavelengthinnanometers}
  \lean{PhyXMiniProblems.ProblemPhyXMini0578.wavelengthInNanometers}
  \uses{def:physics:phyx-mini-0578:phyxminiproblems-problemphyxmini0578-wavelengthquantity, def:physics:phyx-mini-0578:phyxminiproblems-problemphyxmini0578-wavelengthreadout}
  Read a physical wavelength in nanometres
\end{definition}
\begin{proof}
  This is the declared model definition of wavelength In Nanometers.
\end{proof}
\begin{definition}[energy In Joules]
  \label{def:physics:phyx-mini-0578:phyxminiproblems-problemphyxmini0578-energyinjoules}
  \lean{PhyXMiniProblems.ProblemPhyXMini0578.energyInJoules}
  Read a physical energy in coherent-SI joules
\end{definition}
\begin{proof}
  This is the declared model definition of energy In Joules.
\end{proof}
\begin{definition}[planck Constant In Joule Seconds]
  \label{def:physics:phyx-mini-0578:phyxminiproblems-problemphyxmini0578-planckconstantinjouleseconds}
  \lean{PhyXMiniProblems.ProblemPhyXMini0578.planckConstantInJouleSeconds}
  \uses{def:physics:phyx-mini-0578:phyxminiproblems-problemphyxmini0578-planckconstantquantity}
  Read physical action in coherent-SI joule-seconds
\end{definition}
\begin{proof}
  This is the declared model definition of planck Constant In Joule Seconds.
\end{proof}
\begin{definition}[magnetic Field Strength In Teslas]
  \label{def:physics:phyx-mini-0578:phyxminiproblems-problemphyxmini0578-magneticfieldstrengthinteslas}
  \lean{PhyXMiniProblems.ProblemPhyXMini0578.magneticFieldStrengthInTeslas}
  \uses{def:physics:phyx-mini-0578:phyxminiproblems-problemphyxmini0578-magneticfieldstrengthquantity}
  Read magnetic-field strength in coherent-SI teslas
\end{definition}
\begin{proof}
  This is the declared model definition of magnetic Field Strength In Teslas.
\end{proof}
\begin{definition}[magnetic Moment In Joules Per Tesla]
  \label{def:physics:phyx-mini-0578:phyxminiproblems-problemphyxmini0578-magneticmomentinjoulespertesla}
  \lean{PhyXMiniProblems.ProblemPhyXMini0578.magneticMomentInJoulesPerTesla}
  \uses{def:physics:phyx-mini-0578:phyxminiproblems-problemphyxmini0578-magneticmomentmagnitudequantity}
  Read magnetic-moment magnitude in coherent-SI joules per tesla
\end{definition}
\begin{proof}
  This is the declared model definition of magnetic Moment In Joules Per Tesla.
\end{proof}
\begin{definition}[speed Of Light In Meters Per Second]
  \label{def:physics:phyx-mini-0578:phyxminiproblems-problemphyxmini0578-speedoflightinmeterspersecond}
  \lean{PhyXMiniProblems.ProblemPhyXMini0578.speedOfLightInMetersPerSecond}
  Physlib's vacuum light speed, read in metres per second
\end{definition}
\begin{proof}
  This is the declared model definition of speed Of Light In Meters Per Second.
\end{proof}
\begin{definition}[ordinary Planck Constant SIValue]
  \label{def:physics:phyx-mini-0578:phyxminiproblems-problemphyxmini0578-ordinaryplanckconstantsivalue}
  \lean{PhyXMiniProblems.ProblemPhyXMini0578.ordinaryPlanckConstantSIValue}
  Exact coherent-SI calibration of the ordinary Planck constant
\end{definition}
\begin{proof}
  This is the declared model definition of ordinary Planck Constant SIValue.
\end{proof}
\begin{definition}[bohr Magneton SIValue]
  \label{def:physics:phyx-mini-0578:phyxminiproblems-problemphyxmini0578-bohrmagnetonsivalue}
  \lean{PhyXMiniProblems.ProblemPhyXMini0578.bohrMagnetonSIValue}
  Standard coherent-SI calibration of the Bohr magneton
\end{definition}
\begin{proof}
  This is the declared model definition of bohr Magneton SIValue.
\end{proof}
\begin{definition}[Upper Doublet Member]
  \label{def:physics:phyx-mini-0578:phyxminiproblems-problemphyxmini0578-upperdoubletmember}
  \lean{PhyXMiniProblems.ProblemPhyXMini0578.UpperDoubletMember}
  The two members of the spin-split upper 3p level
\end{definition}
\begin{proof}
  This is the declared model definition of Upper Doublet Member.
\end{proof}
\begin{definition}[Atomic Level]
  \label{def:physics:phyx-mini-0578:phyxminiproblems-problemphyxmini0578-atomiclevel}
  \lean{PhyXMiniProblems.ProblemPhyXMini0578.AtomicLevel}
  The common lower 3s level and the two upper 3p levels
\end{definition}
\begin{proof}
  This is the declared model definition of Atomic Level.
\end{proof}
\begin{definition}[upper Atomic Level]
  \label{def:physics:phyx-mini-0578:phyxminiproblems-problemphyxmini0578-upperatomiclevel}
  \lean{PhyXMiniProblems.ProblemPhyXMini0578.upperAtomicLevel}
  \uses{def:physics:phyx-mini-0578:phyxminiproblems-problemphyxmini0578-upperdoubletmember, def:physics:phyx-mini-0578:phyxminiproblems-problemphyxmini0578-atomiclevel}
  Regard an upper-doublet member as its atomic level
\end{definition}
\begin{proof}
  This is the declared model definition of upper Atomic Level.
\end{proof}
\begin{definition}[principal Quantum Number]
  \label{def:physics:phyx-mini-0578:phyxminiproblems-problemphyxmini0578-principalquantumnumber}
  \lean{PhyXMiniProblems.ProblemPhyXMini0578.principalQuantumNumber}
  \uses{def:physics:phyx-mini-0578:phyxminiproblems-problemphyxmini0578-atomiclevel}
  Both pictured manifolds have principal quantum number n = 3
\end{definition}
\begin{proof}
  This is the declared model definition of principal Quantum Number.
\end{proof}
\begin{definition}[orbital Angular Momentum Quantum Number]
  \label{def:physics:phyx-mini-0578:phyxminiproblems-problemphyxmini0578-orbitalangularmomentumquantumnumber}
  \lean{PhyXMiniProblems.ProblemPhyXMini0578.orbitalAngularMomentumQuantumNumber}
  \uses{def:physics:phyx-mini-0578:phyxminiproblems-problemphyxmini0578-atomiclevel}
  The lower level has ℓ = 0; both upper members have ℓ = 1
\end{definition}
\begin{proof}
  This is the declared model definition of orbital Angular Momentum Quantum Number.
\end{proof}
\begin{definition}[Spin Orientation]
  \label{def:physics:phyx-mini-0578:phyxminiproblems-problemphyxmini0578-spinorientation}
  \lean{PhyXMiniProblems.ProblemPhyXMini0578.SpinOrientation}
  The two spin magnetic-moment orientations named in the problem
\end{definition}
\begin{proof}
  This is the declared model definition of Spin Orientation.
\end{proof}
\begin{definition}[zeeman Energy Sign]
  \label{def:physics:phyx-mini-0578:phyxminiproblems-problemphyxmini0578-zeemanenergysign}
  \lean{PhyXMiniProblems.ProblemPhyXMini0578.zeemanEnergySign}
  \uses{def:physics:phyx-mini-0578:phyxminiproblems-problemphyxmini0578-spinorientation}
  Sign in -μ · B: parallel has shift -μB, antiparallel has +μB
\end{definition}
\begin{proof}
  This is the declared model definition of zeeman Energy Sign.
\end{proof}
\begin{definition}[Doublet Line]
  \label{def:physics:phyx-mini-0578:phyxminiproblems-problemphyxmini0578-doubletline}
  \lean{PhyXMiniProblems.ProblemPhyXMini0578.DoubletLine}
  The two sodium-doublet emission lines printed in the source
\end{definition}
\begin{proof}
  This is the declared model definition of Doublet Line.
\end{proof}
\begin{definition}[line Upper Member]
  \label{def:physics:phyx-mini-0578:phyxminiproblems-problemphyxmini0578-lineuppermember}
  \lean{PhyXMiniProblems.ProblemPhyXMini0578.lineUpperMember}
  \uses{def:physics:phyx-mini-0578:phyxminiproblems-problemphyxmini0578-upperdoubletmember, def:physics:phyx-mini-0578:phyxminiproblems-problemphyxmini0578-doubletline}
  The shorter line starts at the higher split upper level
\end{definition}
\begin{proof}
  This is the declared model definition of line Upper Member.
\end{proof}
\begin{definition}[line Lower Level]
  \label{def:physics:phyx-mini-0578:phyxminiproblems-problemphyxmini0578-linelowerlevel}
  \lean{PhyXMiniProblems.ProblemPhyXMini0578.lineLowerLevel}
  \uses{def:physics:phyx-mini-0578:phyxminiproblems-problemphyxmini0578-atomiclevel, def:physics:phyx-mini-0578:phyxminiproblems-problemphyxmini0578-doubletline}
  Both lines end at the same lower 3s level
\end{definition}
\begin{proof}
  This is the declared model definition of line Lower Level.
\end{proof}
\begin{definition}[Figure Level Caption]
  \label{def:physics:phyx-mini-0578:phyxminiproblems-problemphyxmini0578-figurelevelcaption}
  \lean{PhyXMiniProblems.ProblemPhyXMini0578.FigureLevelCaption}
  The two quantum-number captions in image 578
\end{definition}
\begin{proof}
  This is the declared model definition of Figure Level Caption.
\end{proof}
\begin{definition}[Sodium Doublet Figure]
  \label{def:physics:phyx-mini-0578:phyxminiproblems-problemphyxmini0578-sodiumdoubletfigure}
  \lean{PhyXMiniProblems.ProblemPhyXMini0578.SodiumDoubletFigure}
  \uses{def:physics:phyx-mini-0578:phyxminiproblems-problemphyxmini0578-atomiclevel, def:physics:phyx-mini-0578:phyxminiproblems-problemphyxmini0578-doubletline, def:physics:phyx-mini-0578:phyxminiproblems-problemphyxmini0578-figurelevelcaption}
  Literal presentation data from image 578. Display heights and wavelength labels are scalar raster readouts; physical wavelengths are stored separately
\end{definition}
\begin{proof}
  This is the declared model definition of Sodium Doublet Figure.
\end{proof}
\begin{definition}[Atom Species]
  \label{def:physics:phyx-mini-0578:phyxminiproblems-problemphyxmini0578-atomspecies}
  \lean{PhyXMiniProblems.ProblemPhyXMini0578.AtomSpecies}
  Atomic species role named in the problem
\end{definition}
\begin{proof}
  This is the declared model definition of Atom Species.
\end{proof}
\begin{definition}[Sodium Doublet Setup]
  \label{def:physics:phyx-mini-0578:phyxminiproblems-problemphyxmini0578-sodiumdoubletsetup}
  \lean{PhyXMiniProblems.ProblemPhyXMini0578.SodiumDoubletSetup}
  \uses{def:physics:phyx-mini-0578:phyxminiproblems-problemphyxmini0578-wavelengthquantity, def:physics:phyx-mini-0578:phyxminiproblems-problemphyxmini0578-planckconstantquantity, def:physics:phyx-mini-0578:phyxminiproblems-problemphyxmini0578-magneticfieldstrengthquantity, def:physics:phyx-mini-0578:phyxminiproblems-problemphyxmini0578-magneticmomentmagnitudequantity, def:physics:phyx-mini-0578:phyxminiproblems-problemphyxmini0578-upperdoubletmember, def:physics:phyx-mini-0578:phyxminiproblems-problemphyxmini0578-atomiclevel, def:physics:phyx-mini-0578:phyxminiproblems-problemphyxmini0578-spinorientation, def:physics:phyx-mini-0578:phyxminiproblems-problemphyxmini0578-doubletline, def:physics:phyx-mini-0578:phyxminiproblems-problemphyxmini0578-sodiumdoubletfigure, def:physics:phyx-mini-0578:phyxminiproblems-problemphyxmini0578-atomspecies}
  Independent physical data for the sodium doublet. In particular, the internal field is stored rather than defined from wavelengths or an answer choice
\end{definition}
\begin{proof}
  This is the declared model definition of Sodium Doublet Setup.
\end{proof}
\begin{definition}[Matches Excited Sodium Scenario]
  \label{def:physics:phyx-mini-0578:phyxminiproblems-problemphyxmini0578-matchesexcitedsodiumscenario}
  \lean{PhyXMiniProblems.ProblemPhyXMini0578.MatchesExcitedSodiumScenario}
  \uses{def:physics:phyx-mini-0578:phyxminiproblems-problemphyxmini0578-sodiumdoubletsetup}
  The prose describes excited sodium atoms emitting both doublet lines
\end{definition}
\begin{proof}
  This is the declared model definition of Matches Excited Sodium Scenario.
\end{proof}
\begin{definition}[Matches Supplied Sodium Doublet Figure]
  \label{def:physics:phyx-mini-0578:phyxminiproblems-problemphyxmini0578-matchessuppliedsodiumdoubletfigure}
  \lean{PhyXMiniProblems.ProblemPhyXMini0578.MatchesSuppliedSodiumDoubletFigure}
  \uses{def:physics:phyx-mini-0578:phyxminiproblems-problemphyxmini0578-upperatomiclevel, def:physics:phyx-mini-0578:phyxminiproblems-problemphyxmini0578-lineuppermember, def:physics:phyx-mini-0578:phyxminiproblems-problemphyxmini0578-linelowerlevel, def:physics:phyx-mini-0578:phyxminiproblems-problemphyxmini0578-sodiumdoubletsetup}
  Quantum-number labels, bracket, ordering, arrows, and wavelength text visible in the primary image. This contains no magnetic-field value
\end{definition}
\begin{proof}
  This is the declared model definition of Matches Supplied Sodium Doublet Figure.
\end{proof}
\begin{definition}[Matches Sodium Doublet Wavelength Readouts]
  \label{def:physics:phyx-mini-0578:phyxminiproblems-problemphyxmini0578-matchessodiumdoubletwavelengthreadouts}
  \lean{PhyXMiniProblems.ProblemPhyXMini0578.MatchesSodiumDoubletWavelengthReadouts}
  \uses{def:physics:phyx-mini-0578:phyxminiproblems-problemphyxmini0578-wavelengthinnanometers, def:physics:phyx-mini-0578:phyxminiproblems-problemphyxmini0578-sodiumdoubletsetup}
  The printed nanometre labels calibrate the independent physical wavelengths. They are data readouts, not conclusions about the requested field
\end{definition}
\begin{proof}
  This is the declared model definition of Matches Sodium Doublet Wavelength Readouts.
\end{proof}
\begin{definition}[Has Physical Sodium Doublet Quantities]
  \label{def:physics:phyx-mini-0578:phyxminiproblems-problemphyxmini0578-hasphysicalsodiumdoubletquantities}
  \lean{PhyXMiniProblems.ProblemPhyXMini0578.HasPhysicalSodiumDoubletQuantities}
  \uses{def:physics:phyx-mini-0578:phyxminiproblems-problemphyxmini0578-wavelengthinmeters, def:physics:phyx-mini-0578:phyxminiproblems-problemphyxmini0578-energyinjoules, def:physics:phyx-mini-0578:phyxminiproblems-problemphyxmini0578-planckconstantinjouleseconds, def:physics:phyx-mini-0578:phyxminiproblems-problemphyxmini0578-magneticfieldstrengthinteslas, def:physics:phyx-mini-0578:phyxminiproblems-problemphyxmini0578-magneticmomentinjoulespertesla, def:physics:phyx-mini-0578:phyxminiproblems-problemphyxmini0578-upperatomiclevel, def:physics:phyx-mini-0578:phyxminiproblems-problemphyxmini0578-sodiumdoubletsetup}
  Positivity and level ordering required by the physical model
\end{definition}
\begin{proof}
  This is the declared model definition of Has Physical Sodium Doublet Quantities.
\end{proof}
\begin{definition}[Uses Standard Spectroscopic Constants]
  \label{def:physics:phyx-mini-0578:phyxminiproblems-problemphyxmini0578-usesstandardspectroscopicconstants}
  \lean{PhyXMiniProblems.ProblemPhyXMini0578.UsesStandardSpectroscopicConstants}
  \uses{def:physics:phyx-mini-0578:phyxminiproblems-problemphyxmini0578-planckconstantinjouleseconds, def:physics:phyx-mini-0578:phyxminiproblems-problemphyxmini0578-magneticmomentinjoulespertesla, def:physics:phyx-mini-0578:phyxminiproblems-problemphyxmini0578-ordinaryplanckconstantsivalue, def:physics:phyx-mini-0578:phyxminiproblems-problemphyxmini0578-bohrmagnetonsivalue, def:physics:phyx-mini-0578:phyxminiproblems-problemphyxmini0578-sodiumdoubletsetup}
  Standard SI calibrations of ordinary Planck action and the electron spin magnetic-moment magnitude. Neither calibration mentions the requested field
\end{definition}
\begin{proof}
  This is the declared model definition of Uses Standard Spectroscopic Constants.
\end{proof}
\begin{definition}[Satisfies Sodium Radiative Transition Laws]
  \label{def:physics:phyx-mini-0578:phyxminiproblems-problemphyxmini0578-satisfiessodiumradiativetransitionlaws}
  \lean{PhyXMiniProblems.ProblemPhyXMini0578.SatisfiesSodiumRadiativeTransitionLaws}
  \uses{def:physics:phyx-mini-0578:phyxminiproblems-problemphyxmini0578-wavelengthinmeters, def:physics:phyx-mini-0578:phyxminiproblems-problemphyxmini0578-energyinjoules, def:physics:phyx-mini-0578:phyxminiproblems-problemphyxmini0578-planckconstantinjouleseconds, def:physics:phyx-mini-0578:phyxminiproblems-problemphyxmini0578-speedoflightinmeterspersecond, def:physics:phyx-mini-0578:phyxminiproblems-problemphyxmini0578-upperatomiclevel, def:physics:phyx-mini-0578:phyxminiproblems-problemphyxmini0578-lineuppermember, def:physics:phyx-mini-0578:phyxminiproblems-problemphyxmini0578-linelowerlevel, def:physics:phyx-mini-0578:phyxminiproblems-problemphyxmini0578-sodiumdoubletsetup}
  Radiative laws for both lines: each photon carries the atomic energy drop and obeys Eγ = h c / λ. No line is assigned a requested field value
\end{definition}
\begin{proof}
  This is the declared model definition of Satisfies Sodium Radiative Transition Laws.
\end{proof}
\begin{definition}[Satisfies Internal Spin Zeeman Law]
  \label{def:physics:phyx-mini-0578:phyxminiproblems-problemphyxmini0578-satisfiesinternalspinzeemanlaw}
  \lean{PhyXMiniProblems.ProblemPhyXMini0578.SatisfiesInternalSpinZeemanLaw}
  \uses{def:physics:phyx-mini-0578:phyxminiproblems-problemphyxmini0578-energyinjoules, def:physics:phyx-mini-0578:phyxminiproblems-problemphyxmini0578-magneticfieldstrengthinteslas, def:physics:phyx-mini-0578:phyxminiproblems-problemphyxmini0578-magneticmomentinjoulespertesla, def:physics:phyx-mini-0578:phyxminiproblems-problemphyxmini0578-upperatomiclevel, def:physics:phyx-mini-0578:phyxminiproblems-problemphyxmini0578-zeemanenergysign, def:physics:phyx-mini-0578:phyxminiproblems-problemphyxmini0578-sodiumdoubletsetup}
  The general spin Zeeman law E = E₀ + s μ\_B B for each upper member, with sign determined by magnetic-moment orientation. It contains no wavelength or numerical answer for the field
\end{definition}
\begin{proof}
  This is the declared model definition of Satisfies Internal Spin Zeeman Law.
\end{proof}
\begin{definition}[inferred Internal Field In Teslas]
  \label{def:physics:phyx-mini-0578:phyxminiproblems-problemphyxmini0578-inferredinternalfieldinteslas}
  \lean{PhyXMiniProblems.ProblemPhyXMini0578.inferredInternalFieldInTeslas}
  \uses{def:physics:phyx-mini-0578:phyxminiproblems-problemphyxmini0578-wavelengthinmeters, def:physics:phyx-mini-0578:phyxminiproblems-problemphyxmini0578-planckconstantinjouleseconds, def:physics:phyx-mini-0578:phyxminiproblems-problemphyxmini0578-magneticmomentinjoulespertesla, def:physics:phyx-mini-0578:phyxminiproblems-problemphyxmini0578-speedoflightinmeterspersecond, def:physics:phyx-mini-0578:phyxminiproblems-problemphyxmini0578-sodiumdoubletsetup}
  The field inferred from the wavelength readouts and calibrated constants. This scalar prediction does not define the independently stored physical field
\end{definition}
\begin{proof}
  This is the declared model definition of inferred Internal Field In Teslas.
\end{proof}
\begin{definition}[Answer Choice]
  \label{def:physics:phyx-mini-0578:phyxminiproblems-problemphyxmini0578-answerchoice}
  \lean{PhyXMiniProblems.ProblemPhyXMini0578.AnswerChoice}
  Labels of the four field-strength choices supplied with the problem
\end{definition}
\begin{proof}
  This is the declared model definition of Answer Choice.
\end{proof}
\begin{definition}[displayed Field In Teslas]
  \label{def:physics:phyx-mini-0578:phyxminiproblems-problemphyxmini0578-displayedfieldinteslas}
  \lean{PhyXMiniProblems.ProblemPhyXMini0578.displayedFieldInTeslas}
  \uses{def:physics:phyx-mini-0578:phyxminiproblems-problemphyxmini0578-answerchoice}
  Magnetic-field strength, in teslas, printed beside each answer label
\end{definition}
\begin{proof}
  This is the declared model definition of displayed Field In Teslas.
\end{proof}
\begin{definition}[recorded Dataset Answer]
  \label{def:physics:phyx-mini-0578:phyxminiproblems-problemphyxmini0578-recordeddatasetanswer}
  \lean{PhyXMiniProblems.ProblemPhyXMini0578.recordedDatasetAnswer}
  \uses{def:physics:phyx-mini-0578:phyxminiproblems-problemphyxmini0578-answerchoice}
  The answer label recorded by the source dataset, retained as metadata
\end{definition}
\begin{proof}
  This is the declared model definition of recorded Dataset Answer.
\end{proof}
\begin{definition}[Matches Displayed Field]
  \label{def:physics:phyx-mini-0578:phyxminiproblems-problemphyxmini0578-matchesdisplayedfield}
  \lean{PhyXMiniProblems.ProblemPhyXMini0578.MatchesDisplayedField}
  \uses{def:physics:phyx-mini-0578:phyxminiproblems-problemphyxmini0578-magneticfieldstrengthinteslas, def:physics:phyx-mini-0578:phyxminiproblems-problemphyxmini0578-sodiumdoubletsetup, def:physics:phyx-mini-0578:phyxminiproblems-problemphyxmini0578-answerchoice, def:physics:phyx-mini-0578:phyxminiproblems-problemphyxmini0578-displayedfieldinteslas}
  A displayed field agrees when the physical value is within half a tesla of it. The tolerance is a display convention, not a premise fixing the actual field
\end{definition}
\begin{proof}
  This is the declared model definition of Matches Displayed Field.
\end{proof}
\begin{definition}[Is Unique Matching Displayed Field]
  \label{def:physics:phyx-mini-0578:phyxminiproblems-problemphyxmini0578-isuniquematchingdisplayedfield}
  \lean{PhyXMiniProblems.ProblemPhyXMini0578.IsUniqueMatchingDisplayedField}
  \uses{def:physics:phyx-mini-0578:phyxminiproblems-problemphyxmini0578-sodiumdoubletsetup, def:physics:phyx-mini-0578:phyxminiproblems-problemphyxmini0578-answerchoice, def:physics:phyx-mini-0578:phyxminiproblems-problemphyxmini0578-matchesdisplayedfield}
  Exactly one displayed field agrees with the physical field
\end{definition}
\begin{proof}
  This is the declared model definition of Is Unique Matching Displayed Field.
\end{proof}
\begin{lemma}[sodium Doublet Photon Energy Difference]
  \label{lem:physics:phyx-mini-0578:phyxminiproblems-problemphyxmini0578-sodiumdoubletphotonenergydifference}
  \lean{PhyXMiniProblems.ProblemPhyXMini0578.sodiumDoubletPhotonEnergyDifference}
  \uses{def:physics:phyx-mini-0578:phyxminiproblems-problemphyxmini0578-energyinjoules, def:physics:phyx-mini-0578:phyxminiproblems-problemphyxmini0578-magneticfieldstrengthinteslas, def:physics:phyx-mini-0578:phyxminiproblems-problemphyxmini0578-magneticmomentinjoulespertesla, def:physics:phyx-mini-0578:phyxminiproblems-problemphyxmini0578-sodiumdoubletsetup, def:physics:phyx-mini-0578:phyxminiproblems-problemphyxmini0578-matchesexcitedsodiumscenario, def:physics:phyx-mini-0578:phyxminiproblems-problemphyxmini0578-satisfiessodiumradiativetransitionlaws, def:physics:phyx-mini-0578:phyxminiproblems-problemphyxmini0578-satisfiesinternalspinzeemanlaw}
  The shared lower level, energy-drop laws, and Zeeman shifts imply that the photon-energy separation is 2 μ\_B B
\end{lemma}
\begin{proof}
  The conclusion follows from the governing relations and figure readouts named in the statement of sodium Doublet Photon Energy Difference.
\end{proof}
\begin{lemma}[internal Magnetic Field exact Formula]
  \label{lem:physics:phyx-mini-0578:phyxminiproblems-problemphyxmini0578-internalmagneticfield-exactformula}
  \lean{PhyXMiniProblems.ProblemPhyXMini0578.internalMagneticField_exactFormula}
  \uses{def:physics:phyx-mini-0578:phyxminiproblems-problemphyxmini0578-magneticfieldstrengthinteslas, def:physics:phyx-mini-0578:phyxminiproblems-problemphyxmini0578-sodiumdoubletsetup, def:physics:phyx-mini-0578:phyxminiproblems-problemphyxmini0578-matchesexcitedsodiumscenario, def:physics:phyx-mini-0578:phyxminiproblems-problemphyxmini0578-hasphysicalsodiumdoubletquantities, def:physics:phyx-mini-0578:phyxminiproblems-problemphyxmini0578-satisfiessodiumradiativetransitionlaws, def:physics:phyx-mini-0578:phyxminiproblems-problemphyxmini0578-satisfiesinternalspinzeemanlaw, def:physics:phyx-mini-0578:phyxminiproblems-problemphyxmini0578-inferredinternalfieldinteslas}
  Combining the photon relation with the spin splitting gives the exact wavelength formula for the independently stored internal field
\end{lemma}
\begin{proof}
  The conclusion follows from the governing relations and figure readouts named in the statement of internal Magnetic Field exact Formula.
\end{proof}
% --- Archon declaration topology end ---
% --- Archon physics formalization source end ---
